% =====================================================================
% arXiv v2 — FRONT SECTIONS: Abstract + §1 Intro + §2 Related Work + §3 Background
% Frame around the spine (§4–§6). Resolves \label{sec:background}.
% REFRAME 2026-06-11 (Lars): CEP = OPTIONAL, ADVISORY trust convention, not
% binding network-wide enforcement. CEP compels no one; the transaction
% decision always sits with the agent/operator. Scope (b) = "optional
% convention with an operator-independent MATURITY SIGNAL", not a mandatory
% chokepoint. Wording: optional / advisory / convention / demonstrable.
% No "enforces" / "compels" / "provably secure" at the network scope.
% (a)/(b) substance + analogy research + spine mechanics unchanged.
% =====================================================================

\begin{abstract}
Autonomous AI agents now transact at production scale — 69{,}000 bots executing 165 million transactions across \$50 million USDC in cumulative volume on a single marketplace — without any shared trust layer between participants \cite{ref1_agentmarket}. Regulators (Singapore IMDA, NIST CAISI, EU AI Act) and major AI laboratories have converged on the same requirement: an open, verifiable trust infrastructure for agents that no single vendor can deliver alone. MolTrust supplies the verification half — signed Verifiable Credentials and a recomputable Trust Score — while every transaction decision remains with the relying party. The open question is the other half: how a network-wide trust \emph{convention} — a shared default that participating operators voluntarily follow — can become something one relies on without a single party deciding it is in force.

We present the Combined Enforcement Protocol (CEP), framed here as an \textbf{optional, advisory trust convention} on the Agent Authorization Envelope, not a mandatory enforcement mechanism. CEP compels nothing: an operator that follows it gains a network-wide trust signal; one that does not may transact with agents outside the convention — that is the counterparty's choice, not the protocol's. CEP's contribution is a \textbf{publicly recomputable maturity signal} for the convention — a five-condition predicate any party can recompute from anchored data — so that ``the convention is now broadly and diversely adopted'' is a verifiable fact rather than one party's declaration, dissolving the oracle problem such signals usually reintroduce. We distinguish two scopes: \textbf{operator-local} (a relying party fail-closing its own agent against its own mandate — its own authority, buildable today) and \textbf{network-wide} (the optional convention, with its operator-independent maturity signal). The KYC/AML analogy frames how network-wide trust governance historically arises — by convention and distributed consequence, not central coercion — and CEP is one such convention, offered honestly as an approximation, not a solution. The deployed W3C VC + DID infrastructure — live since March 2026 across eight credential verticals, with on-chain anchoring on Base L2 — is the supporting evidence that the substrate is real.
\end{abstract}

% ---------------------------------------------------------------------
\section{Introduction}
\label{sec:intro}
% Material: v1.9 §1 reframed around the convention-maturity question.
% ---------------------------------------------------------------------

The deployment of autonomous AI agents has moved from proof-of-concept into production commerce. A single marketplace reports 69{,}000 autonomous bots executing 165 million transactions across \$50 million USDC in cumulative volume, without any shared trust layer between participants \cite{ref1_agentmarket}. Across financial services, non-human identities now outnumber human employees by wide margins and are growing fast \cite{ref2_neville,ref3_entrolabs}; network data covering a large share of global web traffic shows AI-agent traffic ratios orders of magnitude above human traffic \cite{ref4_prince,ref5_cloudflare}. These are present conditions, not projections.

Regulatory frameworks and industry converge on a structural requirement: an open, portable, cryptographically verifiable trust infrastructure for agents \cite{ref8_imda,ref9_nistcaisi,ref10_nccoe,ref11_anthropic,ref12_google}. W3C Decentralized Identifiers and Verifiable Credentials provide the right primitives \cite{ref6_w3cdid,ref7_w3cvc}, and a growing body of work shows how to assemble them into agent trust systems (Section~\ref{sec:related}). Such infrastructures remove single points of failure for \emph{verification}: any party can check a signed credential without a central service. But verification is only half of trust. The other half is \emph{response} — what a relying party does when an authorization is violated. A relying party can always refuse its own agent's action (Section~\ref{sec:scope}); the harder question is network-wide: how a shared \emph{convention} — a default that participating operators voluntarily follow — becomes something one can rely on, without any single party declaring it in force and thereby recreating the central authority decentralised verification removed. \textbf{Who decides a network-wide trust convention is ready — and how is that not a new central operator?}

This paper's contribution is an answer that does not depend on a person, a chain, or a single instance. We separate two scopes, because they have different bases of authority and the distinction is what makes the question tractable:

\begin{itemize}
  \item \textbf{Operator-local enforcement.} A relying party may gate \emph{its own} agent against \emph{its own} authorization, under its own authority — exactly as any deployer bounds its own software. This needs no convention; it is buildable on the deployed substrate today.
  \item \textbf{Network-wide trust convention.} An \emph{optional} network-wide default that participating operators voluntarily follow — useful only if its readiness (broad, diverse adoption) can be recognised without any single operator's say-so. This is the hard scope, and the one CEP addresses.
\end{itemize}

We are explicit from the outset about what CEP is and is not. It is a \emph{convention}, not a coercive mechanism: the reference layer runs \textbf{advisory} (it verifies and records; it blocks nothing at the network scope), the convention's network-wide adoption is \textbf{voluntary}, and its maturity signal is \textbf{designed and demonstrable on testnet}, not declared in force. The transaction decision always remains with the agent or its operator; CEP offers a signal, never an obligation. A reader expecting a protocol that \emph{enforces} on third parties will not find one here, by design — what we decentralise is the \emph{recognition} that an optional convention has matured, not authority over anyone's agent.

Our contributions are: (1) a precise statement of the convention-maturity problem and its oracle sub-problem — how to recognise an optional network-wide convention as ready without a central authority (Section~\ref{sec:problem}); (2) the Combined Enforcement Protocol, an optional advisory convention whose maturity is certified by a conditions-based, publicly recomputable signal that dissolves the oracle problem (Section~\ref{sec:cep}); (3) an explicit scope-and-limits account, including the operator-local vs.\ network-wide distinction and an honest status (Section~\ref{sec:scope}); and (4) deployment evidence that the verification substrate is real — live since March 2026 across eight verticals, with on-chain anchoring (Section~\ref{sec:evidence}).

% ---------------------------------------------------------------------
\section{Related Work and Positioning}
\label{sec:related}
% Material: v1.9 §2 CONDENSED + extended with the three lineage anchors.
% Western anchors mandatory; never an exclusively non-western source base.
% ---------------------------------------------------------------------

\subsection{Regulatory and institutional context}

The Singapore IMDA \textit{Model AI Governance Framework for Agentic AI} (January 2026) is the first comprehensive governance framework for autonomous agents; it requires each agent to carry a verifiable digital identity and an audit trail of which agent acted under whose authorisation \cite{ref8_imda,ref21_imdafull}. NIST's Center for AI Standards and Innovation launched an AI Agent Standards Initiative (February 2026), and the associated NCCoE concept paper frames the gap directly: agents are commonly treated as generic service accounts without dedicated identity, authorization, or accountability controls \cite{ref9_nistcaisi,ref10_nccoe,ref22_cosais}. The EU AI Act (Regulation (EU) 2024/1689) applies Article~14 (Human Oversight) and Article~15 (Accuracy, Robustness, Cybersecurity) to autonomous agents in high-risk domains \cite{ref23_euaiact}; the allocation of the Article~14 duty is central to our scope position (Section~\ref{subsec:protocol-not-deployer}).

\subsection{Industry framework convergence}

Anthropic's \textit{Trustworthy Agents in Practice} (April 2026) states that the model layer alone cannot secure agentic AI and calls for shared infrastructure no single vendor can provide, favouring open standards over proprietary alternatives \cite{ref11_anthropic}. Google's SAIF~2.0 Agent Risk Map identifies rogue-action and over-permissioned-tool risks and the principles of well-defined controllers, limited powers, and observable actions \cite{ref12_google,ref24_cosai}; Microsoft's Entra Agent ID addresses agent identity within a single enterprise perimeter \cite{ref25_entra}. These efforts establish the need; none answers how a network-wide trust convention reaches relied-upon status across organisational boundaries without a central authority.

\subsection{Academic foundations}

Closest to this work are deployed-or-prototyped agent trust systems. SAGA \cite{ref16_syros} provides centralised, formally analysed policy enforcement through a trusted Provider — strong within one administrative domain, but it presupposes the very trusted party we seek to avoid at the network scope. Hu and Rong \cite{ref15_hurong} give a six-category taxonomy of inter-agent trust (Brief, Claim, Proof, Stake, Reputation, Constraint); Schroeder de Witt \cite{ref14_schroederdewitt} frames multi-agent security as a field with threats specific to agent interaction \cite{ref26_motwani}. Acharya's TIVA \cite{ref17_acharya} and Rodriguez Garzon et al.\ \cite{ref36_rodriguezgarzon} describe DID + VC + on-chain-anchoring architectures conceptually; Rodriguez Garzon et al.\ explicitly identify the limitation of an LLM in sole charge of security procedures, motivating trust mechanisms that operate independently of the model layer. Threat taxonomies by Ferrag et al.\ \cite{ref18_ferrag} (cross-mapped to CVE/NVD) and Mao et al.\ \cite{ref20_mao} catalogue the attack surface, alongside the OWASP agentic-security material \cite{ref19_owasp}. None of these treats the \emph{governance of an optional network-wide trust convention} — how its readiness is recognised without a central authority — as the object of design.

\subsection{Governance, data-availability, and Sybil-resistance lineages}
\label{subsec:lineages}

CEP draws three lineages from outside the agent-trust literature, each of which informs a specific mechanism in Section~\ref{sec:cep}.

\textbf{On-chain governance and concentration.} Analyses of DAO voting power show that on-chain governance frequently concentrates in a few actors, with low Nakamoto coefficients and weak participation \cite{ref39_daogov}. This is the baseline CEP's concentration caps (Section~\ref{subsec:params}) are designed to beat, and the comparison point for why timelocks and per-actor/per-cluster bounds matter.

\textbf{Optimistic-rollup security (verification \texttt{>} production).} Optimistic rollups establish correctness not by trusting a block producer but by letting any verifier recompute and challenge a published claim \cite{ref38_arbitrum}. CEP's honest-verifier data availability (Section~\ref{subsec:hvda}) borrows this stance directly: the convention's maturity predicate is a deterministic function anyone can recompute from anchored data, so no party holds the oracle.

\textbf{Sybil resistance from graph structure.} Social-graph Sybil defences \cite{ref40_sybilguard} and proof-of-personhood / unique-humanity systems \cite{ref41_brightid} establish that resistance to identity fabrication is better grounded in structural relationships than in self-declared attributes. CEP's cluster-diversity condition (Section~\ref{subsec:clusters}) follows this lineage: independence is computed from the endorsement graph, not from declared categories.

\subsection{Position summary}

This paper is complementary to the verification-and-identity work above and orthogonal to it in one respect: its object is not another way to verify or to enforce, but the \emph{recognition} that an optional network-wide trust convention has matured — a signal no single party should own. That object is what scope (b) and CEP address; the deployed substrate (Section~\ref{sec:background}) is the evidence that the question is not hypothetical.

% ---------------------------------------------------------------------
\section{Background: An Advisory Verification Substrate}
\label{sec:background}
% Material: v1.9 §3.1–3.4 CONDENSED + §17.1 advisory framing. Short; feeds spine.
% ---------------------------------------------------------------------

This section condenses the deployed substrate to just enough to ground the governance argument; full normative definitions are in the Protocol Technical Specification \cite{ref28_techspec}. The point to carry forward is that the substrate exists and runs \textbf{advisory} — it verifies and records, it does not block at the network scope — so the open question is not how to force enforcement on it, but how an optional network-wide convention built on it becomes something operators can rely on.

\textbf{Four primitives.} Each agent holds a W3C Decentralized Identifier \cite{ref6_w3cdid}, with ownership proven by an Ed25519 key; the reference method is \texttt{did:moltrust}, but any method supporting Ed25519/secp256k1 is conformant. Authorization is expressed as a W3C Verifiable Credential \cite{ref7_w3cvc} supporting attenuating delegation chains any verifier can traverse independently. Behaviour is recorded as Interaction Proof Records — dual-signed, Merkle-anchored on Base L2 — accumulating into a registry-computed Trust Score (Section~\ref{sec:evidence}). All three are expressed in open, portable formats, so evidence issued by one party is verifiable by any other.

\textbf{The Agent Authorization Envelope (AAE).} The machine-evaluable authorization object carries three normative blocks: \textbf{MANDATE} (permitted purpose, allowed/denied action patterns, delegation rules), \textbf{CONSTRAINTS} (time windows, financial thresholds, jurisdictions, counterparty minimums, obligations), and \textbf{VALIDITY} (issuer, holder binding, mandatory expiry, revocation, optional anchor). Any conformant evaluator applies default-deny, deny-precedence, attenuation-only delegation, and mandatory expiry. The design reuses NIST SP~800-162 ABAC \cite{ref29_abac}, RFC~9396 \cite{ref30_rar}, RFC~8785 JCS \cite{ref31_jcs}, and Ed25519Signature2020 \cite{ref32_ed25519}.

\textbf{Three-layer architecture, advisory today.} Authorization is checked at three layers: \emph{cryptographic} (Ed25519 + RFC~8785, tamper-evident by construction), \emph{API} (trust-score and revocation consequences, fail-closed on unreachable revocation), and \emph{kernel} (Falco eBPF detection below the agent's process boundary \cite{ref13_falcobridge}). In the reference deployment these run \textbf{advisory}: a DENY is produced as a signed, anchorable verdict and recorded, but the action is not blocked at the network scope, and whether to act on the verdict is the relying party's decision \cite{ref28_techspec}. Operator-local blocking (scope (a), Section~\ref{subsec:protocol-not-deployer}) is available to a relying party in its own gateway; at the network scope nothing is blocked — the convention of Section~\ref{sec:cep} is something operators opt into, not a chokepoint imposed on them.

\textbf{Decentralised evidence, centralised registry — for now.} The evidence layer is genuinely verifier-independent: any party with an issuer's key validates any signed artifact without a central service. The reference registry, however, is operated by one organisation, and how an optional network-wide convention on that substrate would be recognised as mature — without that recognition resting on the single operating organisation — is, today, an open design question. That question is the subject of the rest of the paper, and the problem it poses is the subject of the next section.
